
\appendix

\subsection{Contributions}

CX, LK started the project. CX built the infrastructure for online RL. JTS designed and trained the RL token. ME built the intervention interface. AA and AE designed and built the gripper and robot hardware. CX, LK designed the system implementation, task suite and the experiments. SL, LK provided advice throughout the project. LK, CX, JTS, SL, ME worked on writing, illustrations and the video. 

\subsection{Additional experiment details}
\label{app:exp}

% \kk{@Charles can you write this section, as if you are explaining to Michael how to reproduce our exp on Ethernet}

First, we collect a demonstration dataset on the target task; we then fine-tune the base VLA model and train the RL token for 2000 to 10000 gradient steps on the single task data. The VLA is then frozen during online RL training.

During online RL, we initialize the RL actor and critic from scratch with a two-layer MLP (hidden dimension 256) for the zip tie fastening, Ethernet, and charger insertion tasks. We use a larger network consisting of a three-layer MLP with hidden dimension 512 for the more challenging screw installation task. Both networks receive the RL token produced by the frozen base VLA model, proprioceptive position and velocity as input. The critic is trained with an ensemble of two Q functions following \citet{fujimoto2018td3} and we use the minimum of the two Q functions to calculate target values. The actor additionally takes in the reference action chunk produced by the VLA model, which is masked out at 50\% during training and always provided during inference. The actor is parameterized as a Gaussian policy with a small fixed standard deviation that outputs an action chunk $\mathbf{a}_{t:t+C-1} \in \mathbb{R}^{C \times d}$ with $C{=}10$ from the current observation. To increase sample efficiency, we subsample action chunks 2 control steps apart during training, so each second of data produces roughly 25 samples for the RL network. A sparse +1 reward is provided by the operator during training when the RL task has been completed. 

For the screw installation and zip tie fastening tasks, we first start the RL training in the critical-phase setting only. We then advance to a full task phase, first running the base model to complete the non-critical phase of the task, and switching to the RL policy when the critical phase is reached. This two-phase training strategy improves training efficiency while ensuring the RL policy is robust to the initial distributions induced by the base policy on the earlier part of the task. We report the policy performance after gathering about 5 hours of data. 


\subsection{Additional experimental details for the baselines}
For all baseline methods, we use the same environment and action space setup as for our method -- policies execute in delta action space at 50\,Hz.
\label{app:baseline}

\textbf{PLD}: Following the original paper, we first pre-train the critic network on 50 base policy rollouts with Cal-QL~\cite{nakamoto2023cal} for better sample efficiency. We then proceed to the online RL stage.

\textbf{DSRL}: Following the original implementation, our implementation predicts a $(1, 32)$ dimensional latent action, which is repeated $50$ times in the first dimension to match the noise input space of our action-chunk VLA.

\textbf{HIL-SERL}: Following the original implementation, we initialize RLPD training with 20 episodes of demonstrations and provide interventions throughout training. However, it was not able to succeed in our setting because of the higher control frequency (50\,Hz) compared to the original system (10\,Hz) and the absence of an action-space bounding box to reduce exploration space. 

\textbf{DAgger}: We fine-tune our VLA with a mixture of demonstration data and the same set of intervention data collected during the online RL training. 